\section{An arbitrary-gap birth payment and a finite owner theorem}
The source-gap constants five and six are sharp for pointwise
nonrecovery on their respective owner-count ranges. They do not
prevent an unbounded theorem in which the allowable skeleton
degree varies with the source gap. Here such a theorem is proved
without a finite computation or a fixed upper bound on the gap.

Let $Z=1^\alpha U$ and $Y$ be actual increasing arms, with
$h=|U|\ge2$, $g\ge2$, and write
\[
 a=\operatorname{wt}Z,\quad D=\sum_{t\in U}(t-1)=a-\alpha-h,
 \quad\kappa=\kappa_Y,\quad d=\kappa-\alpha.
\]
At shared forward rank $j\ge2$, suppose
\[
 F_Z(j+1)\le F_Z(j),\qquad
 E_Y(0),\ldots,E_Y(j+1)\ \hbox{is nondecreasing}.       \tag{1}
\]
These are the same mixed/source-prefix premises as before.
Put $n=a+\operatorname{wt}Y-j$ and
$T=e_{n+1}-e_n$ for $e=(2-z)q_Zq_Y$.

\begin{theorem}[Finite owner at every first missing weight]
If $T>0$ under (1), then
\[
 \boxed{d\ge6,\qquad D\le2d(d-3),\qquad
       a\le\alpha+4d(d-3),\qquad
       j\le\alpha+4d(d-3)+1.}                         \tag{2}
\]
In particular, for each fixed $\kappa$, there are only finitely
many possible owner arms and shared ranks with positive recovery.
Explicit bounds, independent of every opposite capacity and count, are
\[
 \boxed{a\le1+4(\kappa-1)(\kappa-4),\qquad
        j\le2+4(\kappa-1)(\kappa-4).}                 \tag{3}
\]
Equivalently, if $d\ge3$ and $D>2d(d-3)$, the mixed tail
cannot recover. With the usual nonnegative OLD A2 head input,
this gives S at every head, without C, a half-weight premise,
or a restriction on the first missing weight.
\end{theorem}

\begin{proof}
Use the exact source identity and primitive notation from the
preceding source-birth theorem:
\[
 T=(s_{r+1}-s_r)-\Delta(E_YF_Z)_{j+1},\quad
 r=a-j,\quad s_r=\Delta q_Z(r),\quad
 \Delta(E_YF_Z)_{j+1}\ge F_Z(j+1-\kappa)               \tag{4}
\]
whenever the displayed birth index is present. All terms of the
source difference convolution through this rank are nonnegative.
The proved uniform gap-five theorem already excludes $d\le5$;
thus assume $d\ge6$ when proving necessity in (2).
For the stated sufficient large-$D$ version, $3\le d\le5$
is already covered by that same earlier theorem.

If $r\le-2$, the primitive curvature is zero and $T\le0$.
If $r=-1$, the curvature is one. Unless $D+h+2-d<0$,
the birth coefficient in (4) lies on $F_Z$'s support and is
at least one. Hence a positive recovery in this case requires
$D\le d-h-3\le d-5\le2d(d-3)$, as required.

For $r\ge0$, a positive recovery forces
$d_{r+1}>d_r>0$ for the ordinary positive primitive
$d_r=(\Delta q_Z(r))^+$. Set
\[
       e=D-2r-2\ge0,\qquad C=d_{r+1}-d_r>0.
\]
The skeleton-only relative envelope and degree rank ratio give
\[
                         (r+1)C\le(e+1)d_r.           \tag{5}
\]
If $s=d-e-2\le0$, the first birth is present, and lower-center
reflection gives $F_Z(j+1-\kappa)\ge q_Z(r+1)\ge C$.
Thus necessarily $s>0$. Write $L=s-1=d-e-3\ge0$.

We claim that
\[
                    r\ge dL+e                         \tag{6}
\]
is already sufficient to pay the entire curvature by the FIRST
birth. This is a simpler, gap-uniform sufficient condition for
the earlier finite-product criterion. It applies at every gap.
Condition (6) implies $r\ge L$, so all following indices exist.
For $0\le i<L$ put
\[
             x_i=\frac{e+3+2i}{r-i}\ge0,
             S=\sum_{i=0}^{L-1}x_i,\qquad
             P=\prod_{i=0}^{L-1}(1+x_i).
\]
The sharper primitive relative envelope gives the backward chain
\[
 d_{r-i}\le(1+x_i)d_{r-i-1},\qquad d_r\le P d_{r-L}.   \tag{7}
\]
Let $B=r-L+1>0$. Since $B\le r-i$ at each step,
\[
 BS\le\sum_{i=0}^{L-1}(e+3+2i)
      =L(e+L+2)=L(d-1).
\]
Consequently
\[
 B(1-S)\ge r+1-dL\ge e+1>0,
 \quad S<1,\quad (r+1)(1-S)\ge e+1.                 \tag{8}
\]
For arbitrary nonnegative $x_i$ of sum at most one,
\[
                     (1-S)\prod_i(1+x_i)\le1.         \tag{9}
\]
Here is the complete elementary induction: after a partial sum
$u$ and product $V\ge0$, insertion of $x\ge0$ multiplies the
budget by a coefficient satisfying
$(1-u-x)(1+x)\le1-u$, because the difference is $x(u+x)$.
Every partial sum is at most one, so the induction applies.
Empty sums and products are included.
Combining (8)--(9) gives $(e+1)P\le r+1$.
Together with (5)--(7), this proves
\[
                              C\le d_{r-L}.           \tag{10}
\]
The reflected pair
$u=r-L=r-s+1$, $v=r+\kappa-1$ has endpoint sum
$u+v=\alpha+D=a-h$. Both indices are nonnegative,
and $u\le v$. The established lower-center comparison therefore
gives
\[
 F_Z(j+1-\kappa)=q_Z(v)\ge q_Z(u)\ge d_u\ge C.
\]
The birth index is present: explicitly it equals
$r+h-L\ge h$. Equations (4) and (10) contradict $T>0$.
This proves the claim (6).

Any positive recovery with $r\ge0$ must therefore satisfy
\[
 r<d(d-e-3)+e,
 \quad r+(d-1)e<d(d-3).
\]
All parameters are integers. Hence
\[
 D=2r+2+e\le2d(d-3)-(2d-3)e\le2d(d-3).             \tag{11}
\]
Since $h\le D$, this proves (2), including $j\le a+1$.
For fixed $\kappa$, $\alpha=\kappa-d\ge1$ and
$6\le d\le\kappa-1$. The expression
$\kappa-d+4d(d-3)$ increases over this integer interval:
its difference from the endpoint value is
$(\kappa-1-d)(4(\kappa-1+d)-13)\ge0$.
Thus (3) follows. There are finitely many $\alpha,h,D$ and
partitions of $D$ into $h$ positive capacity increments, so
these are genuine finite-owner conclusions, not just rank bounds.
\end{proof}

Every step of the arbitrary-length product budget, backward
chain, ratio-sum estimate and scalar payment, together with the
integer owner/weight implications, is checked in
\texttt{ArbitrarySourceGapFiniteOwner.lean}. These are unbounded
generic results. Their identification with the previously proved
actual-arm skeleton and source identities is the written proof
above. No end-to-end Lean composition datatype is asserted.

\section{A count-only corollary of the exact head-mass criterion}
For $K=\lfloor(2g+1)/5\rfloor+1$, the earlier mass polynomial satisfies
\[
 A(g,L_Y)=[z^K](1+z)^{g-L_Y+4}[3]^{L_Y}
           \ge {g+4\choose K}\ge2^K.                 \tag{12}
\]
Indeed $2K\le g+4$ for every $g\ge2$; choose $K$ disjoint
pairs of the $g+4$ binary variables and exactly one variable
from each pair. Their $2^K$ distinct subsets all have rank $K$.
If $2g\ge5j+4$, then $K\ge j+2$ and
$A(g,L_Y)\ge2^{j+2}\ge5\,2^{j-1}$.
Thus the earlier general mixed mass theorem proves S at ALL
heads in this count range, without OLD A2 or any source-prefix
premise. For a mixed S edge still outside that theorem,
\[
                               2g<5j+4.               \tag{13}
\]
Combining (3) and (13), all still-uncovered mixed S edges with
a fixed opposite $\kappa$ and a nondecreasing opposite SOURCE
prefix have bounded owner weight, shared forward rank, and
opposite non-one count. Opposite capacities and padding may
remain unbounded; no complete finite reduction for every S
guard is claimed here.

These conclusions do not settle both-forward-rising late source
overlap or the mixed region after the opposite SOURCE prefix
has declined. They also do not prove S for the remaining finite
owner ranges at each $\kappa$, and therefore do not imply general
C/A2 or original unamplified coefficient unimodality. The main
conjecture remains OPEN. The twelve scoped kernel reports and
independent exact diagnostics are kept distinct from a proof
of that conjecture.

The subsequent theorem in \texttt{source\_curvature\_linear\_owner.tex}
uses genuine later source increments to strengthen the degree bound
to $D\le12d-25$ and the fixed-first-weight owner bound to
$a\le24\kappa-73$. The quadratic bounds above remain valid and
are stronger for the smallest source gaps.
